\input{../stock/en-tete_v6.tex}

\newcommand{\titrechapitre}{titre}
\newcommand{\numerochapitre}{numéro}

\renewcommand{\headrulewidth}{0.4pt}
\renewcommand{\footrulewidth}{0.4pt}
\fancyfoot[L]{Notes de Raphaël JONTEF}
\fancyfoot[C]{\thepage}
\fancyfoot[R]{Enseirb-Matmeca}
\fancyhead[L]{Cours de Mathématiques}
\fancyhead[R]{\titrechapitre}

\begin{document}
\input{../stock/commands.tex}

% Page de titre custom
\setlength{\headheight}{15pt}
\begin{titlepage}
    \centering
    \vspace*{3cm}
    {\huge Chapitre \numerochapitre\par}
    {\Huge \bfseries\titrechapitre\par}
    \vspace{2cm}
    {\Large Notes rédigées par Raphaël JONTEF\par}
    \vspace{0.5cm}
    {\normalsize Cours enseigné par François DUFOUR\par}
    \vspace{4cm}
    {\small Enseirb-Matmeca\par}
    {\small filière informatique\par}
    \vfill
    {\large \today\par}
\end{titlepage}

% plan du cours
\tableofcontents
\newpage


% -----------Corps du document--------------

\begin{proposition}{première}{un équivalent}
    On a~:
    $$\frac{(2n)!}{(n!)^2} \symbolelim{n \to + \infty}{}{\sim} \frac{4^n}{\sqrt{\pi n}}$$
\end{proposition}

\begin{demonstration}
    On applique la formule de Stirling~:
    \begin{align*}
        (2n)! &\symbolelim{n \to + \infty}{}{\sim} \sqrt{4 \pi n} \left(\frac{2n}{e}\right)^{2n} \\
        (n!)^2 &\symbolelim{n \to + \infty}{}{\sim} 2 \pi n \left(\frac{n}{e}\right)^{2n}
    \end{align*}
    On en déduit~:
    $$\frac{(2n)!}{(n!)^2} \symbolelim{n \to + \infty}{}{\sim} \frac{4^n}{\sqrt{\pi n}}$$
\end{demonstration}

\begin{proposition}{finale}{équivalent de la quantité en jeu}
    On a~:
    $$\binom{2n-1}{n-1} \symbolelim{n \to + \infty}{}{\sim} \frac{1}{2\sqrt{\pi}} \frac{4^n}{\sqrt{n}}$$
\end{proposition}

\begin{demonstration}
    On a~:
    \begin{align*}
        \binom{2n-1}{n-1} &= \frac{(2n-1)!}{n!(n-1)!}\\
        &= \frac{n}{2n} \frac{(2n)!}{(n!)^2} \\
        &\symbolelim{n \to + \infty}{}{\sim} \frac{1}{2} \frac{4^n}{\sqrt{\pi n}} \\
    \end{align*}
    D'où le résultat et le caractère exponentiel en $n$ de cette quantité.
\end{demonstration}

\end{document}